\section{Rule-Deviation Decoding}
\label{sec:method}

% 证据：rulegen/ruledec.py、neural/ruledec_ext.py（排班解码）；bench/fjsp2.py、fjsp_base3.py（FJSP：sp 均匀备选、spr 规则组合备选）；
%       evalrule.py（代理分）；neural/gp.py、nnrule.py、es.py、train_il.py；rulegen.py（LLM）；实验5续文档（ρ 的选取）。
% 2026-10-27 按新框架重写：加入备选动作集 A(s)、与优化器无关；学习器细节移到补充材料 S-IV。

\subsection{Overview}
\label{sec:method:overview}

\RDD divides the work between a construction rule and the genotype.
The rule makes every decision by default.
The genotype does not encode every decision.
It encodes a few \emph{deviations}: the decisions at which the decoder should not follow the rule, and what it should do there instead.
Fig.~\ref{fig:framework} shows how \RDD is used.
A rule is learned offline on training instances (Section~\ref{sec:method:learn}).
An MOEA then searches over the deviations, and every genotype it creates is decoded relative to the rule.

\begin{figure*}[t]
  \centering
  \figplaceholder{5.2cm}{Fig.~2 (framework, to be drawn): (1) offline rule learning with the proxy score; (2) online search with the RDD decoder and the deviation rate $\rho$; (3) applications: static rostering, re-rostering, flexible job shop.}
  \caption{Overview of rule-deviation decoding. A rule is learned offline by scoring the solutions that random deviations from it produce. An MOEA then searches over deviation slots; the decoder follows the rule wherever no slot is active.}
  \label{fig:framework}
\end{figure*}

\subsection{Rules and Alternatives}
\label{sec:method:rule}

A schedule is built one decision at a time.
At decision $i$ the decoder is in state $s_i$, and the rule $\rul$ returns its first choice $\rul(s_i)$.
A deviation replaces this choice by an action from an \emph{alternative list} $A(s_i)$.
The alternative list is part of the design.
We use two kinds:
\begin{itemize}
  \item \textbf{Ranked alternatives.} $A(s_i)$ holds all other feasible actions in the rule's order.
  This suits problems with few feasible actions per decision.
  We use it in rostering, where a deviation either sends one of the few other available robots or opens a new robot.
  \item \textbf{Rule portfolio.} A few other rules each propose their first choice, and $A(s_i)$ holds these proposals.
  A proposal may coincide with the rule's choice; the deviation then has no effect.
  This suits problems with many feasible actions, most of them poor.
  We use it in the flexible job shop, where a decision chooses among dozens of operation--machine pairs and many standard dispatching rules exist.
\end{itemize}
In both cases a deviation picks one entry of the list, and under random genes every entry is equally likely.
What matters is that the entries are reasonable candidates.
On the flexible job shop, a rule portfolio is clearly better than ranked alternatives over all operation--machine pairs (Section~\ref{sec:exp-ablation}).

In rostering, the rule sees the current time, the available robots, and for every robot its duty count, hours, remaining budget, duties left before maintenance and ready time.
It also sees the parameters of each robot type and a \emph{gene type} $\tau_g$, a robot type read from one gene, which the rule may use when it opens a robot.
The simplest rule is the \emph{default rule} $\rul_0$: send the available robot with the fewest duties, and open a robot of type $\tau_g$ if none is available.

\subsection{Genotype and Decoding}
\label{sec:method:genotype}

The genotype holds $\K$ \emph{deviation slots}, two genes per slot, in $[0,1)$:
\begin{equation}
  \genome = (p_1, a_1, \dots, p_\K, a_\K, \text{other genes}).
\end{equation}
In rostering, the other genes are the gene type $g$ and one tool gene per shift, read as in the baselines; FJSP needs no other genes.
The number of slots grows with the number $N$ of decisions, $\K = \max(4, \lceil \dens N \rceil)$, where $\dens$ is the \emph{slot ratio}.
Each slot works as follows:
\begin{itemize}
  \item \textbf{On or off.} The slot is active if $a_k \ge 1 - \qon$, where $\qon$ is the \emph{activation probability}.
  \item \textbf{Where.} An active slot acts on decision $\lfloor p_k N \rfloor$. If two active slots hit the same decision, the later one counts.
  \item \textbf{What.} The value $u_k = (a_k - 1 + \qon)/\qon \in [0,1)$ picks an alternative from $A(s)$ by its position in the set.
\end{itemize}
Algorithm~\ref{alg:decode} lists the decoder.
It collects the active slots, then builds the schedule decision by decision.
At a decision with an active slot it takes the chosen alternative, if that action is feasible; otherwise it takes the rule's first choice.
In rostering the alternatives are split as in the baselines: values $u_k < 0.85$ send another robot and larger values open a new robot of a type chosen by $u_k$;
the decoder also keeps the capacity check and the backtracking repair of \dec{5} on decisions without a deviation.

\begin{algorithm}[t]
  \caption{Rule-deviation decoding}
  \label{alg:decode}
  \begin{algorithmic}[1]
    \Require genotype $\genome$, rule $\rul$, alternative lists $A(\cdot)$, slot ratio $\dens$, activation probability $\qon$
    \Ensure a schedule, or \textsc{infeasible}
    \State $\K \gets \max(4, \lceil \dens N \rceil)$;\ \ $\Delta \gets \emptyset$
    \For{$k = 1, \dots, \K$}
      \If{$a_k \ge 1 - \qon$} \State $\Delta[\lfloor p_k N \rfloor] \gets (a_k - 1 + \qon)/\qon$ \EndIf
    \EndFor
    \State $s \gets$ initial state
    \For{$i = 0, 1, \dots$ until the schedule is complete}
      \State $\act \gets \rul(s)$
      \If{$i \in \Delta$}
        \State $\act' \gets$ element of $A(s)$ at position $\lfloor \Delta[i]\, |A(s)| \rfloor$
        \If{$\act'$ is feasible in $s$} $\act \gets \act'$ \EndIf
      \EndIf
      \State $s \gets$ apply $\act$ to $s$
      \If{a resource limit is exceeded} \Return \textsc{infeasible} \EndIf
    \EndFor
    \State \Return the schedule in $s$
  \end{algorithmic}
\end{algorithm}

The key difference from per-decision encodings lies in what happens after a change.
In a per-decision encoding, a changed decision changes the state, and every later gene is then read in this new state, often with a new meaning.
In \RDD, most later decisions have no gene.
The rule makes them, and it adapts to the new state.
Fig.~\ref{fig:encoding} illustrates this on ten shifts.

\begin{figure}[t]
  \centering
  \begin{tikzpicture}[
      sh/.style={draw, minimum width=5.6mm, minimum height=5mm, inner sep=0pt, font=\scriptsize},
      gn/.style={font=\scriptsize},
      sl/.style={draw, rounded corners=1pt, minimum height=4.5mm, inner sep=1.5pt, font=\scriptsize},
      ar/.style={-{Latex[length=1.5mm]}, thin}]
    \node[font=\footnotesize, anchor=west] at (-0.4,0.75) {(a) Per-decision encoding (\dec{5})};
    \foreach \i in {1,...,10} {
      \node[sh] (a\i) at (0.64*\i,0) {\i};
      \node[gn] (x\i) at (0.64*\i,-0.75) {$x_{\i}$};
      \draw[ar] (x\i) -- (a\i);
    }
    \node[font=\footnotesize, anchor=west] at (-0.4,-1.55) {(b) Rule-deviation decoding (\RDD)};
    \foreach \i in {1,2,4,5,6,8,9,10} {
      \node[sh, fill=black!10] (b\i) at (0.64*\i,-2.3) {\i};
    }
    \node[sh, very thick] (b3) at (0.64*3,-2.3) {3};
    \node[sh, very thick] (b7) at (0.64*7,-2.3) {7};
    \node[sl] (s1) at (1.3,-3.45) {slot 1: shift 3, 2nd choice};
    \node[sl] (s2) at (4.15,-3.45) {slot 2: shift 7, open};
    \node[sl, dashed, text=black!50] (s3) at (6.05,-3.45) {slot 3: off};
    \draw[ar] (s1.north) -- (b3.south);
    \draw[ar] (s2.north) -- (b7.south);
    \node[font=\scriptsize, anchor=west] at (6.7,-2.3) {rule $\rul$};
  \end{tikzpicture}
  \caption{Two encodings of a ten-shift roster. (a) Each shift has its own gene. (b) Gray shifts follow the rule $\rul$; the genotype holds deviation slots, here two active ones that change shifts 3 and 7.}
  \label{fig:encoding}
\end{figure}

\subsection{The Deviation Rate}
\label{sec:method:rate}

In a uniformly random genotype, each slot is active with probability $\qon$, so about $\qon\K \approx \qon\dens N$ slots are active.
We call
\begin{equation}
  \drate = \qon\dens
\end{equation}
the \emph{deviation rate}, the expected number of active slots per decision.
It sets the balance between bias and freedom.
A small $\drate$ keeps random genotypes close to the rule's schedule; a large $\drate$ lets them spread far from it.
On both problems, settings with the same $\drate$ but different $\dens$ and $\qon$ perform alike (Section~\ref{sec:experiments}), so $\drate$ is the parameter to set.
The realized share of decisions that differ from the rule is smaller than $\drate$, because several slots may hit the same decision and an alternative may coincide with the rule's choice.

Two existing encodings are limiting cases of \RDD.
A per-decision encoding that ranks the options by a rule corresponds to one gene for every decision and no default: every decision is free.
Heuristic-space search \cite{storer1992new,dorndorf1995evolution}, which picks a rule for every decision, corresponds to a rule portfolio with no default rule.
\RDD adds what both lack: a default, and a parameter that sets how often the search may leave it.

\textbf{Optimizer.}
\RDD is an encoding, so any optimizer for real-valued genotypes can use it.
Our main experiments use NSGA-II \cite{deb2002nsga2}.
We also report MOEA/D \cite{zhang2007moead} and a Pareto local search that, starting from the rule's schedule, repeatedly redraws $j$ random slots of an archived solution.

\subsection{Learning the Rule}
\label{sec:method:learn}

In \RDD the rule is used inside a search.
A good rule is therefore not simply one that builds a good schedule on its own.
It is one near which the search finds good schedules.
We measure this with a \emph{proxy score}.
For each training instance $I$ we fix a set $G_I$ of 67 genotypes: three with all slots off and 64 uniformly random ones.
Let $\mathcal{D}_\rul(\genome)$ be the solution that \RDD decodes from $\genome$ with rule $\rul$. The proxy score is
\begin{equation}
  \proxy(\rul) = \frac{1}{|\mathcal{I}_{\text{train}}|} \sum_{I \in \mathcal{I}_{\text{train}}} \HV_I\bigl(\{\mathcal{D}_\rul(\genome) : \genome \in G_I\}\bigr),
\end{equation}
where $\HV_I$ is the hypervolume of the non-dominated feasible solutions, with the objectives normalized per instance by the solutions of the default rule $\rul_0$.
All rules are scored on the same genotypes, so differences between rules are not due to sampling.
One evaluation of $\proxy$ costs $15 \times 67 = 1005$ decodings.
During learning we use $\drate = 0.1$.

The proxy score treats the rule as a black box, so any kind of rule can be optimized with it.
We use three learners; their settings are given in Section~S-IV of the supplementary material.
\begin{itemize}
  \item \textbf{GP.} Two trees score \reuse and \new candidates from 22 and 17 size-independent features of the state and the candidate (population 64, 30 generations).
  \item \textbf{Neural network + ES (ES16).} A network with 16 hidden units scores each candidate from 27 features.
  It is first trained to imitate good rosters found by \dec{5}, which gives a useful starting point, and then OpenAI-ES \cite{salimans2017evolution} optimizes $\proxy$ directly.
  \item \textbf{LLM.} An LLM writes the rule as a Python function. We evolve it in the style of EoH \cite{liu2024evolution} for 12 generations, showing the LLM the score of each rule; one run needs about 55 LLM calls.
\end{itemize}
For FJSP we do not learn a rule; we use a standard dispatching rule.

\subsection{Setting the Deviation Rate}
\label{sec:method:k}

The best deviation rate depends on how far good solutions lie from the rule and on the budget (Section~\ref{sec:theory}).
Neither is known for a new instance, so we set $\drate$ once on training instances, with a selection procedure fixed before the runs.
For rostering, each candidate setting was run with the eight GP and ES16 rules on the 15 training instances,
and the setting with the highest mean HV was kept.
We first compared $\drate = 0.1$ to $0.4$ with $\qon = 0.4$, and then $\qon = 0.4$, 0.6 and 0.8 at the largest slot ratio.
The procedure selected $\drate = 0.4$ ($\dens = 1.0$, $\qon = 0.4$).\footnote{Before this selection, an exploratory run had evaluated $\drate = 0.2$ on the test instances. The final rate was chosen on the training instances only; test results at other rates are reported only as a sensitivity analysis.}
For FJSP, the settings of \RDD and of all compared methods were chosen in the same way on 15 instances that are not used in the comparison (Section~\ref{sec:exp-fjsp}).
